\SourceLeaf{392}%
\Needspace{9\baselineskip}
\begin{center}\bfseries Bemerkungen.\end{center}
\noindent\hspace*{1.2em} Am 11. December 1822 hatte \textsc{Gauss} als Beantwortung einer von der Königlichen Societät der Wissenschaften in Copenhagen aufgegebenen Preisfrage seine Abhandlung \textit{Allgemeine Auflösung der Aufgabe: die Theile einer gegebenen Fläche auf einer andern gegebenen Fläche so abzubilden, dass die Abbildung dem Abgebildeten in den kleinsten Theilen ähnlich wird} (Bd. IV. S. 189 bis 216) an \textsc{Schumacher} abgesandt. Die vorstehenden Aufzeichnungen vom 13. bis 15. December 1822, die auf einem einzelnen Blatte gemacht sind, beweisen, dass er ihr mit voller Absicht \textsc{Newtons} stolzes Motto »Ab his via sternitur ad maiora« auf den Weg gegeben hatte, denn die Darstellung des Quadrates des Linienelementes einer beliebigen Fläche in der Form
\[
mm(dt^2+du^2)
\]
hatte ihm den Zugang zu der Lösung des Problems eröffnet, das Krümmungsmass, dessen Unveränderlichkeit gegenüber Biegungen er schon früh erkannt hatte (vergl. S. 372), durch die Coefficienten des Linienelementes, also hier durch $mm$, und die Ableitungen davon auszudrücken, und damit war die Grundlage für die \textit{Disquisitiones generales circa superficies curvas} gewonnen.

Auf denselben Gegenstand bezieht sich eine Notiz in einem Handbuche, die vermuthlich ungefähr aus derselben Zeit, wie die bereits S. 373 abgedruckte Aufzeichnung stammt

{\itshape »Wahrscheinlich findet folgendes höchst zierliche allgemeine Theorem statt:

Ist $n=e^\nu$, so wird
\[
\begin{aligned}
\frac{\partial\partial\nu}{\partial x^2}+\frac{\partial\partial\nu}{\partial y^2}
&=\frac{n\left(\frac{\partial\partial n}{\partial x^2}+\frac{\partial\partial n}{\partial y^2}\right)-\left\{\left(\frac{\partial n}{\partial x}\right)^2+\left(\frac{\partial n}{\partial y}\right)^2\right\}}{nn}\\
&=\frac1{rr'nn}=\frac{d\Sigma}{dS.nn}=\frac{d\Sigma}{ds},
\end{aligned}
\]
wenn $r$ und $r'$ die beiden äussersten Krümmungshalbmesser, die auf einander senkrecht,

\noindent [bedeuten und] wenn
\begin{center}
\begin{tabular}{ccccc}
$d\Sigma$&Raum&auf&der&Himmelskugel\\
$dS$&,,&,,&,,&krummen Fläche\\
$ds$&,,&\multicolumn{3}{l}{in der Ebene}
\end{tabular}
\end{center}
[und $ds=nn\,dS$ ist.]«\par}
\begin{flushright}\textsc{Stäckel.}\end{flushright}
